% ======================================================================
% sections/conclusions.tex
%
% Conclusions scaffold for the paper:
%   Triple Humanoid 24/7 Adverse Event Oncology Trial Response Team:
%   4-Site Rotation
%
% This file is a draft scaffold. A future Claude Code Opus 4.7 1M Max
% session reads the bracketed instruction blocks below, then reads the
% files those instructions name, then expands the scaffold into the
% finished conclusions section. Target length is roughly 3 to 5 pages
% of the 70 plus page final paper.
% ======================================================================

\begin{sectionaccent}

[CONCLUSIONS INSTRUCTIONS FOR FUTURE CLAUDE CODE OPUS 4.7 1M MAX
SESSION. Write 4 short subsections inside this single
\texttt{\textbackslash section} block. Conclusions does not introduce new
files. Every file named here must also be cited earlier in the paper.

Subsection 1, Thesis Restatement. Restate the thesis verbatim. Do not
paraphrase. The thesis appears in
\texttt{demo-projects/07-humanoid/paper/instructions/README.md} under
the heading ``Thesis''.

Subsection 2, Evidence Roll Up. List the 10 evidence anchors that
support the thesis. Each anchor is one bullet, one sentence:
\begin{itemize}
  \item \texttt{codegen/src/swarm\_coordinator.py} (191 lines)
        implements the camaraderie logic that animates the 3 robot
        Unitree H2 EDU camarade swarm at one site. The compact
        line count is itself evidence of clean coordination logic.
  \item \texttt{codegen/src/fda\_rtct\_submitter.py} (67 lines)
        submits CTCAE grade 2 or higher records to the FDA RTCT
        pilot intake API within 1 hour per the HR 9505 SLA, using
        ed25519 signing in production.
  \item The 27 files in \texttt{codegen/src/} span Python, C++,
        Rust, YAML, and Markdown and together implement the full
        sensor to LLM to robot pipeline.
  \item \texttt{codegen/data/iterations/index.jsonl} carries 32
        deterministic iteration sweeps across the 4 trial sites and
        many columns.
  \item \texttt{codegen/reports/dashboard.html} adds project
        interpretability for future paper readers.
  \item \texttt{codegen/src/check\_errors.py} provides LLM
        credibility in generating error checking code.
  \item \texttt{codegen/tests/test\_xyz\_safety.py} provides LLM
        credibility in generating test code.
  \item \texttt{codegen/src/h2\_dispatcher.py} provides per-site H2
        dispatch proficiency.
  \item \texttt{execution/data/site\_runtime\_runs/AT-01\_llm\_decisions.jsonl}
        plus the matching SF-01, SD-01, and BO-01 files show 1 Hz
        decision logs at scale for all 4 sites.
  \item \texttt{execution/data/iterations\_aggregate.duckdb} and the
        JSONL files in \texttt{execution/data/} together
        demonstrate DuckDB and JSONL proficiency in the same data
        layer.
\end{itemize}

Subsection 3, Industry Checklist Framing. Present the entire body of
work as a sequence of checklist items that the industry no longer
has to redo. Pull the checklist from the Inventory at a Glance table
in
\texttt{demo-projects/07-humanoid/paper/instructions/README.md}. Mark
each item as done so the next team can build on top of it. The point
of the framing is to save the industry the cost of repeating the
groundwork.

Subsection 4, Forward Statement. Close with the United States
patient safety, efficacy, and speed framing. The United States needs
to remain Number 1 in the world regarding patient safety, efficacy,
and speed benefits in oncological adverse event robot swarms for
clinical trials. Frame the closing sentence as an invitation to the
FDA, sponsors, sites, and IRBs to advance the new field of Physical
AI oncology trial surgical procedures together. Cite the author's
December 2025 paper
\texttt{kawchak\_2025\_18029100} one final time as the prior
evidence that LLMs are effective on FDA adverse event data.

Required real-life feasibility insight to include in this section:
``The pipeline in this paper is small enough to audit, large enough
to deploy at 4 sites, and tight enough to honor the FDA RTCT 1 hour
SLA. The work is ready for the human governance step.''

Do not exceed 5 pages. Use single dashes only. Use \S\ for the
section sign. Avoid stranding 1 or 2 word lines on their own.]

The conclusions section will restate the thesis verbatim, roll up the
10 evidence anchors, frame the body of work as an industry checklist
of completed steps, and close with the United States patient safety,
efficacy, and speed forward statement plus the final citation to
\texttt{kawchak\_2025\_18029100}.

\sectiontable{Conclusions planning matrix}{%
Thesis & C1 & 1 & instructions README & Verbatim restatement & 1 paragraph & No paraphrase & Ready \\
Evidence & C2 & 2 & 10 anchors & Bullet roll up & 1 page & 10 bullets & Ready \\
Checklist & C3 & 3 & Inventory at a Glance & Done checklist framing & 1 to 2 pages & Industry handoff & Ready \\
Forward & C4 & 4 & kawchak\_2025\_18029100 & US leadership close & 1 paragraph & Final citation & Ready \\
}

\end{sectionaccent}
